\documentclass[12pt,a4paper]{article}
\usepackage[a4paper,margin=2.2cm]{geometry}
\usepackage{fontspec}
\usepackage{polyglossia}
\setdefaultlanguage{arabic}
\setotherlanguage{english}
\newfontfamily\arabicfont[Script=Arabic,Scale=1.15]{Amiri}
\newfontfamily\englishfont{TeX Gyre Termes}
\usepackage{enumitem}
\usepackage{setspace}
\usepackage{array}
\setlength{\parindent}{0pt}
\setlength{\parskip}{7pt}
\setstretch{1.25}
\begin{document}

\begin{flushright}
\textbf{بروف بابكر م. عثمان}\\
القاهرة\\
جامعة النهضة\\
واتساب: 00249965428114\\
الهاتف المحمول: 00201031628478\\
البريد الإلكتروني: \textenglish{babikerosman@yahoo.com}
\end{flushright}

\begin{flushleft}
٢٦ أغسطس ٢٠٢٦
\end{flushleft}

\vspace{1em}
\begin{center}
\Large\textbf{إلى من يهمه الأمر}
\end{center}

وُلد مهند في المملكة العربية السعودية عام 2006، وأصيب بالتهاب الدماغ الفيروسي عندما كان عمره سنة ونصف تقريبًا، وذلك أثناء زيارة لقضاء العطلة في السودان.

وللأسف، لم يتم تشخيص حالته في البداية بصورة فورية ودقيقة. وبعد ذلك عُرض على عيادة الأطفال في السودان، حيث شُخّصت حالته على أنها شلل دماغي ناقص التوتر العضلي.

أظهر التصوير بالرنين المغناطيسي للدماغ ضمورًا دماغيًا منتشرًا وتوسعًا في الخزانة المخية الكبيرة (\textenglish{Cisterna Magna}) دون وجود تغيرات بؤرية أخرى.

وأظهر تصوير بالرنين المغناطيسي لاحق وجود ضمور دماغي مصحوب بتلين بيضاء حول البطينات سابق لأوانه، مع درجة من ضمور المخيخ.

تشير هذه الخصائص إلى وجود حالة اعتلال وعائي دقيق (\textenglish{Microangiopathic state}) يُرجح أن تكون ناجمة عن نقص أكسجة سابق أو إصابة وعائية، وهو ما يتوافق مع المرض السابق والتشنجات المطولة غير المعالجة التي حدثت في الأبيض.

في المملكة العربية السعودية، تم تشخيص مهند عادةً بمتلازمة لينّوكس--غاستو (\textenglish{Lennox--Gastaut Syndrome -- LGS}) وهي حالة تتميز بوجود أنواع متعددة من نوبات الصرع، وتأخر النمو، وضعف القدرات المعرفية، ومقاومة للعلاجات التقليدية.

أما أحدث تشخيص، والذي وضعه أطباء أعصاب الأطفال السعوديون، فهو:

\begin{center}
\textbf{الشلل الدماغي، أعقبه اعتلال دماغي نخري حاد في مرحلة الطفولة}\\
\textenglish{Acute Necrotizing Encephalopathy of Childhood}
\end{center}

حاليًا، لا يستطيع مهند الجلوس أو التدحرج إلى جانبه، رغم قدرته على تحريك ذراعيه وساقيه، مع كون أداء الجانب الأيمن أفضل بصورة ملحوظة من الجانب الأيسر.

ويتلقى تغذيته باستخدام حليب \textenglish{F-100} عبر أنبوب تغذية معدي. وتشمل أدويته مضادات الصرع والفيتامينات والمضادات الحيوية الوقائية، على النحو التالي:

\begin{enumerate}[label=\arabic*. ,leftmargin=*,itemsep=3pt]
\item باراسيتامول (400 ملغ ثلاث مرات يوميًا).
\item إيبوبروفين (200 ملغ ثلاث مرات يوميًا).
\item سيبتازول (قرص واحد مرتين يوميًا).
\item فلاجيل (500 ملغ مرتين يوميًا).
\item فينوباربيتال (50 ملغ مرتين يوميًا).
\item كلونازيبام (25.0 ملغ ثلاث مرات يوميًا).
\item ليفيتيراسيتام (900 ملغ مرتين يوميًا).
\item لاموترجين (50 ملغ مرتين يوميًا).
\item دومبيريدون (10 ملغ مرتين يوميًا).
\item إيزوميبرازول (20 ملغ مرتين يوميًا).
\item ألفيرين سبازم (ألفيرين + ديميثيكون) (قرص واحد مرتين يوميًا).
\item إنتيروجيرمينا (بروبيوتيك) (أمبولة واحدة مرتين يوميًا).
\item \textenglish{Enemax} (مرة كل يومين).
\item كولي يورينال (خلين + هيكسامين + بيبرازين) (مرتين يوميًا).
\item حليب تركيبة \textenglish{F-100} (350 مل مرتين يوميًا).
\item أكياس محلول الإماهة الفموية \textenglish{ORS} (كيس واحد عند الحاجة).
\item شراب متعدد الفيتامينات (\textenglish{Sansovit} مع الحديد) (15 مل مرة واحدة يوميًا).
\item الكالسيوم (شراب \textenglish{Osteocare}) (10 مل مرة واحدة يوميًا).
\end{enumerate}

\vspace{1em}
مع خالص التحية والتقدير،

\textbf{بروف بابكر م. عثمان}\\
أكاديمي وباحث في العلوم الطبية الحيوية

\end{document}
